\paragraph{R2-W1, R2-D1 (Index creation w.r.t.\ $|\mathcal{W}|$ and update overheads).}

\begin{itemize}[leftmargin=1.2em, itemsep=1pt, topsep=2pt]
\item \textit{W1: ``more analysis and discussion around the index creation/update overheads.''}
\item \textit{D1: ``The constraint-aware index creation (w.r.t.\ $|\mathcal{W}|$) and update overheads discussion could be interesting.''}
\end{itemize}
\smallskip

\noindent\textbf{Response.}
We add an explicit cost analysis. As stated in Section~2.2, the index design follows the structure of $\omega$: metric-aware structures (\emph{``weighted KD-trees, ball trees, cover trees''})~\cite{reif2022scalable} for dimension-reweighting $\omega$, and \emph{``partitioned indices, bitmap filtering, or inverted indices''}~\cite{lee2019constraint,chan1998bitmap,babenko2014inverted} for discrete-constraint $\omega$. Rather than one index per specification, $\mathcal{W}$ is partitioned into a small number of structural families that share an index, so construction cost scales with the number of families, not $|\mathcal{W}|$, and is amortized across queries. Our scalability study (Table~\ref{tab:scale}) shows the online set-selection cost is essentially constant in $N$, so retrieval and index maintenance dominate cost, motivating the constraint-aware index.

\noindent\emph{Revision (Appendix~A, ``Practical realization and costs''):}
\begin{quote}
\textcolor{blue}{The offline index is amortized across queries: the specification family $\mathcal W$ is partitioned into a small number of structural families (metric-reweighting vs.\ discrete-constraint dominated), each served by one reusable index, so construction cost scales with the number of families rather than $|\mathcal W|$. Under insertions and deletions to $\mathcal D$ the index is maintained incrementally, and a query whose $\omega\notin\mathcal W$ falls back to a linear scan, so the cost model degrades gracefully.}
\end{quote}
